% Figure: schematic Fanno line (adiabatic flow with friction) and
% Rayleigh line (frictionless flow with heat addition) on a
% temperature-entropy diagram. Both lines have a sonic point (M=1) at
% their right extremity; the upper branch is subsonic, the lower
% supersonic. A normal shock jumps from the supersonic point on one
% curve to the subsonic point at the same mass flux and momentum.
\begin{figure}[htbp]
\centering
\begin{tikzpicture}
\begin{axis}[
  width=12cm, height=8cm,
  xlabel={entropy $s$ (arbitrary units)}, ylabel={temperature $T$ (arbitrary units)},
  xmin=0, xmax=10, ymin=0, ymax=10,
  axis lines=left,
  grid=major, grid style={enggray!20},
  tick label style={font=\scriptsize},
  label style={font=\small},
  legend style={font=\scriptsize, at={(0.03,0.97)}, anchor=north west},
  legend cell align=left,
]
% Fanno line: a curve that rises, turns at a sonic point (max entropy),
% and falls. Parametrise with a simple lobe. Use plot coordinates.
\addplot[engblue, very thick, smooth] coordinates {
  (1.0,8.6) (2.2,8.3) (3.6,7.8) (5.2,7.0) (6.6,6.0)
  (7.8,4.6) (8.6,3.0) (8.95,1.9) (8.6,1.1) (7.9,0.6)};
\addlegendentry{Fanno line (friction, adiabatic)}
% Rayleigh line: a fuller loop; max entropy at sonic, max temperature
% at a subsonic point left of sonic.
\addplot[engred, very thick, smooth] coordinates {
  (1.0,5.0) (2.6,6.2) (4.2,7.1) (5.8,7.5) (7.0,7.3)
  (8.0,6.5) (8.7,5.2) (9.0,3.6) (8.7,2.2) (8.0,1.2)};
\addlegendentry{Rayleigh line (heat addition, frictionless)}
% sonic point markers near the right tips (max entropy)
\addplot[only marks, mark=*, engblack, mark size=1.6pt] coordinates {(8.95,1.9)};
\addplot[only marks, mark=square*, engblack, mark size=1.6pt] coordinates {(9.0,3.6)};
\node[engblack, font=\scriptsize, anchor=west] at (axis cs:9.05,1.9) {$M=1$};
\node[engblack, font=\scriptsize, anchor=west] at (axis cs:9.1,3.6) {$M=1$};
% branch labels
\node[engblue, font=\scriptsize] at (axis cs:3.0,8.7) {subsonic};
\node[engblue, font=\scriptsize] at (axis cs:7.6,0.6) {supersonic};
\end{axis}
\end{tikzpicture}
\caption[Fanno and Rayleigh lines on a $T$-$s$ diagram]{Schematic Fanno
and Rayleigh lines on a temperature-entropy diagram, the two
one-dimensional flows that drive a constant-area duct toward $M=1$ from
either side. The Fanno line traces adiabatic flow with wall friction at
fixed mass flux: friction always raises entropy, so a subsonic flow on
the upper branch accelerates toward the sonic point at the right tip,
and a supersonic flow on the lower branch decelerates toward the same
point; the duct chokes when the flow reaches $M=1$ and added length
cannot pass more flow. The Rayleigh line traces frictionless flow with
heat addition at fixed mass flux and momentum: heating likewise drives
both branches toward the sonic point, so a combustor can thermally
choke. Each curve has its maximum entropy at $M=1$, which is why a
flow cannot cross from subsonic to supersonic by friction or heating
alone. The construction is the working tool for duct flow with friction
(Fanno) and with heat release (Rayleigh), the topics of the worked
examples below.}
\label{fig:vol03ch13:fanno-rayleigh}
\end{figure}
